% !TeX root = ../../Book5.tex
% 纲要标题：### 9.3 区间表示与线性定向箭图
% 纲要位置：工作记录/计划纲要.md，第 4023 行。

\section{区间表示与线性定向箭图}
\label{b5:sec:0903}


% 纲要范围：
%
% * A_2、A_3 和线性定向 A_n
% * 直和与不可分解表示
% * 用线性代数分解区间表示
% * 区间同态的方向条件与显式计算
%

一个箭头的分类只需矩阵的秩。沿一条有向直线增加箭头以后，需要同时记录每段复合的秩；这些数仍然足以恢复分解。本节给出实际选基的办法，为后面更一般的箭图分类提供可计算的比较。

\subsection{\texorpdfstring{\(A_2\)、\(A_3\) 与线性定向 \(A_n\)}{A2、A3 与线性定向 An}}
\label{b5:subsec:0903:01}

\begin{definition}\label{b5:def:interval-quiver-representation}
	设 \(k\) 为域，\(n\ge1\) 为整数，\(Q\) 为线性定向箭图 \(1\to2\to\cdots\to n\)。对满足 \(1\le a\le b\le n\) 的整数 \(a,b\)，在顶点 \(a,a+1,\ldots,b\) 放置 \(k\)，在其余顶点放置零空间；区间内部的箭映射为恒等，其他箭映射为零。所得表示记为 \(I[a,b]\)，称为\term[qujian-biaoshi]{区间表示}[interval representation]。
\end{definition}
\symidx{\(I[a,b]\)：线性定向箭图的区间表示}
对 \(A_2=(1\to2)\)，给定映射 \(f:V_1\to V_2\)，令 \(r=\operatorname{rank}f\)、\(d_i=\dim V_i\)。在 \(V_1\) 中补全 \(\ker f\) 的基，并用其余基向量的像作为 \(\operatorname{im}f\) 的基，再补全到 \(V_2\)，得到
\[
	V\cong I[1,2]^{\oplus r}\oplus I[1,1]^{\oplus(d_1-r)} \oplus I[2,2]^{\oplus(d_2-r)}.
\]
因此两端空间的维数与箭的秩共同决定同构类；只给维数则不够。
\ProseParagraphBreak
对于 \(1\xrightarrow{f}2\xrightarrow{g}3\)，还须比较 \(\operatorname{im}f\) 与 \(\ker g\)。先在它们的交中选基，分别补到两个子空间，再补到 \(V_2\)；把 \(\operatorname{im}f\) 的这些基向量提升到 \(V_1\)，补上 \(\ker f\) 的基，并在 \(V_3\) 中补全相应的像。这就把表示分成六种区间。记 \(r_f,r_g,r_{gf}\) 为三个映射的秩，则相应重数为
\[
	\begin{aligned}
		m_{13}&=r_{gf},&m_{12}&=r_f-r_{gf},&m_{23}&=r_g-r_{gf},\\
		m_{11}&=d_1-r_f,&m_{22}&=d_2-r_f-r_g+r_{gf},& m_{33}&=d_3-r_g.
	\end{aligned}
\]
重数 \(m_{22}\ge0\) 来自
\[
	\dim(\operatorname{im}f\cap\ker g)=r_f-r_{gf} \le\dim\ker g=d_2-r_g.
\]
这个不等式正好给出所需的非负性；对应的向量是 \(\ker g\) 中不来自前一顶点的部分。
\ProseParagraphBreak
另一种定向 \(1\xrightarrow{f}2\xleftarrow{g}3\) 也可直接分解。记 \(d_i=\dim_kV_i\)、\(r_f=\operatorname{rank}f\)、\(r_g=\operatorname{rank}g\)，并令
\[
	[f\ g]:V_1\oplus V_3\longrightarrow V_2,\qquad (x,y)\longmapsto f(x)+g(y),\qquad r=\operatorname{rank}[f\ g].
\]
对两支箭的像 \(U=\operatorname{im}f\)、\(W=\operatorname{im}g\)，有 \(\operatorname{im}[f\ g]=U+W\)，所以
\[
	\dim_k(U\cap W)=r_f+r_g-r.
\]
取交空间的基，分别补到 \(U,W\)，再补到 \(V_2\)，并把 \(U,W\) 的基分别提升到两端。交中的每个基方向给一个在三个顶点都为 \(k\)、两箭均为恒等的表示，记为 \(I_{123}\)；\(U,W\) 各自的补给相邻两顶点的表示 \(I_{12},I_{23}\)，其中支集内的箭也是恒等；其余顶点空间为零。剩余基及两端的核给单点表示，记为 \(S_1,S_2,S_3\)。六种模型的重数因此为
\[
	\begin{aligned}
		m_{123}&=r_f+r_g-r,&m_{12}&=r-r_g,&m_{23}&=r-r_f,\\
		m_1&=d_1-r_f,&m_2&=d_2-r,&m_3&=d_3-r_g.
	\end{aligned}
\]
前三个数分别是 \(U\cap W\)、\(U\) 中所选补空间、\(W\) 中所选补空间的维数，后三个数分别来自 \(\ker f\)、\(V_2/(U+W)\)、\(\ker g\)，因而均非负。这些模型的自同态也只有标量，且维数向量互不相同，故六个数决定分解及同构类。在线性定向中，复合秩测量从前一顶点进入的向量有多少继续传播；在叉形定向中，像的交测量从两端进入的方向有多少重合。

\subsection{直和与不可分解表示}
\label{b5:subsec:0903:02}

\begin{definition}\label{b5:def:indecomposable-quiver-representation}
	设 \(k\) 为域，\(Q\) 为有限箭图，\(V\) 为 \(Q\) 的非零有限维 \(k\) 表示。若 \(V\) 不能写成两个非零表示的直和，就称其为\term[jiantu-bukefenjie]{不可分解表示}[indecomposable representation]。若 \(V\) 没有非零真子表示，就称其为简单表示。两者不同：在箭图 \(1\to2\) 中，区间表示 \(I[1,2]\) 不可分解，却有支集只在顶点2的子表示。
\end{definition}
一个区间表示的自同态在每个非零顶点都是标量，恒等箭迫使这些标量相同。因此
\[
	\operatorname{End}_Q\bigl(I[a,b]\bigr)\cong k.
\]
若它有非平凡直和分解，相应投影会给出不等于零、一的幂等自同态，与域中无此幂等元矛盾。所以每个区间不可分解。不同区间的维数向量不同，故不同构。

\ProseParagraphBreak

不可分解表示的自同态未必都可逆。取 \(A=k[t]/(t^2)\)、左 \(A\)-模 \(X=A\)，也就是单圈箭图加关系 \(t^2=0\) 的二维表示。自同态由一的像唯一确定，故 \(\operatorname{End}_A(X)\cong A\)，各自同态都是乘以 \(a+bt\)。当 \(a\ne0\) 时，逆元为 \(a^{-1}-a^{-2}bt\)；当 \(a=0\) 时，\((bt)^2=0\)。非单位元恰组成理想 \((t)\)，这个环只有零、一两个幂等元，所以 \(X\) 不可分解。这里可以看到局部自同态环的一项作用：非单位元的有限和仍是非单位元。

\ProseParagraphBreak

下面使用 Krull--Schmidt 分解\footnote{克鲁尔（Wolfgang Krull，1899-1971），德国数学家，研究可换代数与模的分解；施密特（\OriginalName{Отто Юльевич Шмидт}，1891-1956），苏联数学家，也研究地球物理与天文学。}。
\begin{proposition}[Krull--Schmidt]\label{b5:thm:quiver-krull-schmidt}
	设 \(k\) 为域，\(Q\) 为有限箭图，\(I\) 为 \(kQ\) 的双边理想。每个满足关系 \(I\) 的有限维 \(k\) 表示都能分解为有限个不可分解表示的直和；各直和项的同构类型及重数唯一到重排。无关系情形取 \(I=0\)。
\end{proposition}
\begin{proof}
	对给定表示 \(V\)，每次把一个非零可分解直和项 \(X\) 拆成 \(Y\oplus Z\)，其中 \(Y,Z\) 都非零，非零直和项的个数便增加一。各项的总维数之和始终为 \(\sum_{i\in Q_0}\dim_kV_i\)，而每个非零项至少占一维，所以项数不超过这个数。因此拆分过程有限步终止，得到存在性；零表示对应空直和。 为证唯一性，先看不可分解表示 \(X\) 的自同态 \(u\)。核与像的升降链在有限维下稳定，取充分大的 \(N\)，便有
	\[
		X=\ker u^N\oplus\operatorname{im}u^N.
	\]
	确实，可取 \(N\) 使 \(\ker u^{2N}=\ker u^N\) 且 \(\operatorname{im}u^{N+1}=\operatorname{im}u^N\)。若 \(x=u^Ny\) 又满足 \(u^Nx=0\)，则 \(u^{2N}y=0\)，稳定性给出 \(u^Ny=0\)，所以交为零；秩与零化度之和为 \(\dim X\)，便得到直和。核与像都是子表示，而 \(X\) 不可分解，故其中一项为零。若核为零，\(u\) 单射，有限维性使它可逆；若像为零，则 \(u^N=0\)。因此每个自同态都或者可逆，或者幂零。

	其自同态环的非单位元构成理想：左右乘任意自同态不可能把不可逆有限维线性算子变可逆；若两个非单位元 \(u,v\) 的和可逆，则 \(a=(u+v)^{-1}u\) 与 \(1-a=(u+v)^{-1}v\) 都非单位，否则再乘 \(u+v\) 就使 \(u\) 或 \(v\) 可逆。对自同态 \(a\) 应用刚证的结论，得 \(a^r=0\) 对某个 \(r\) 成立；但 \(1+a+\cdots+a^{r-1}\) 是 \(1-a\) 的逆，矛盾。 现在设 \(V=\bigoplus_iX_i=\bigoplus_jY_j\)。把 \(X_1\) 包含到 \(V\)，再投到 \(Y_j\)，得到 \(\alpha_j:X_1\to Y_j\)；把 \(Y_j\) 包含到 \(V\)，再投回 \(X_1\)，得到 \(\beta_j:Y_j\to X_1\)。各 \(Y_j\) 上的坐标投影之和为 \(1_V\)，故
	\[
		1_{X_1}=\sum_j\beta_j\alpha_j.
	\]
	若每个复合都不可逆，它们的有限和也属于非单位元理想，便不可能等于恒等映射。因此至少一个 \(\beta_j\alpha_j\) 可逆，\((\beta_j\alpha_j)^{-1}\beta_j\) 给出 \(\alpha_j\) 的左逆。相应映射 \(X_1\to Y_j\) 为分裂单射；\(Y_j\) 不可分解，故它是同构。更具体地，\(V\to X_1\) 在这个 \(Y_j\) 上为同构，所以 \(V=Y_j\oplus\bigoplus_{i>1}X_i\)：任意元可先减去唯一具有相同 \(X_1\) 投影的 \(Y_j\) 元。与原来的 \(Y_j\) 分解取商，得到 \(\bigoplus_{i>1}X_i\cong\bigoplus_{\ell\ne j}Y_\ell\)。 按直和项数归纳即得唯一性。
\end{proof}
此论证只使用有限维性，没有使用路径代数本身有限维。因此带圈的有限箭图，只要所讨论的表示有限维，也满足这个分解唯一性结论。

\subsection{区间表示的线性代数分解}
\label{b5:subsec:0903:03}

\begin{theorem}\label{b5:thm:linear-quiver-intervals}
	设 \(k\) 为域，\(n\ge1\) 为整数，\(Q=(1\to2\to\cdots\to n)\)，\(V=(V_i,V_a)\) 为有限维 \(k\) 表示。则 \(V\) 是区间表示 \(I[a,b]\) 的直和，其中 \(1\le a\le b\le n\)。记
	\[
		r_{ij}=\operatorname{rank}_k(V_i\to V_j)\quad(1\le i\le j\le n), \qquad r_{ii}=\dim_k V_i,
	\]
	其中 \(V_i\to V_j\) 为沿唯一有向路径的箭映射复合，并约定 \(r_{0j}=r_{i,n+1}=0\)，则 \(I[a,b]\) 的重数为
	\[
		m_{ab}=r_{ab}-r_{a-1,b}-r_{a,b+1}+r_{a-1,b+1}.
	\]
	特别地，恰有 \(n(n+1)/2\) 个不可分解同构类。
\end{theorem}
\begin{proof}
	对 \(n\) 归纳。\(n=1\) 是选基。假定已把前 \(n-1\) 个顶点分成区间。在顶点 \(n-1\) 已为零的区间不参与最后一支箭，可直接保留。其余各区间都结束于 \(n-1\)，把其末端空间记为 \(E=V_{n-1}\)。按这些区间的起点定义旗
	\[
		0=F_0\subseteq F_1\subseteq\cdots\subseteq F_{n-1}=E,
	\]
	其中 \(F_i\) 由起点不超过 \(i\) 的末端基向量生成，亦即 \(V_i\to E\) 的像。令 \(f:E\to V_n\) 为最后的箭，\(K=\ker f\)。可以选一组同时适应旗 \(F_i\) 与子空间 \(K\) 的基。逐步把 \(K\cap F_{i-1}\) 的已选基补全为 \(K\cap F_i\) 的基，同时把商空间 \((F_{i-1}+K)/K\) 的已选基补全为 \((F_i+K)/K\) 的基，将新商基向量提升到 \(F_i\)，先前的提升保持不变。这些提升的像在商中线性无关，所以它们与所选的 \(K\cap F_i\) 基合在一起仍无关；任一 \(F_i\) 中元素先用提升消去商类，剩余就在 \(K\cap F_i\)，故这组向量生成 \(F_i\)。到最后一步便得到 \(E\) 的基，其中一部分又是 \(K\) 的基。 任意保持这面旗的可逆换基都延伸为这些末端区间之和的自同构。为看清这一点，将起点为 \(a\) 的末端向量送到起点 \(b\le a\) 的末端向量，给出 \(I[a,n-1]\to I[b,n-1]\) 的同态：在重合区间取同一个标量，在其余顶点取零。旗保持的矩阵恰由这些允许的项组成；其逆也保持旗，故同样延拓，且两复合在每个顶点都是恒等。 例如，\(I[1,3]\oplus I[3,3]\) 在顶点 \(3\) 的末端基记为 \(e,f\)，其中 \(e\) 来自长区间。前两个顶点传播到末端的像都是 \(ke\)，故旗为 \(0\subseteq ke\subseteq ke\oplus kf\)。换成 \(e,f+e\) 保持这面旗：只在顶点 \(3\) 出现的向量可以加上早先传播来的向量。反之，把 \(e\) 送到 \(f\) 不保持 \(ke\)，因而不能与顶点 \(2\to3\) 的箭交换。现在接上末箭
	\[
		h:ke\oplus kf\longrightarrow k,\qquad h(e)=1,\quad h(f)=-1.
	\]
	它的核是 \(k(e+f)\)，在旧基中并非某个区间的末端坐标。取 \(e'=e\)、\(f'=e+f\)，既保持原来的传播像 \(ke'=ke\)，又有 \(h(e')=1\)、\(h(f')=0\)。于是长区间延长到顶点 \(4\)，新选出的短区间仍结束于顶点 \(3\)，得到 \(I[1,4]\oplus I[3,3]\)。这一换基同时使核成为坐标子空间，并保持此前的箭映射。 因此可以用刚才选出的适应基代替原区间末端基。属于 \(K\) 的向量使相应区间仍结束于 \(n-1\)；其余向量的像线性无关，将它们作为 \(\operatorname{im}f\) 的基，相应区间便延长到 \(n\)。最后补全到 \(V_n\) 的基，新增的向量各给一个 \(I[n,n]\)。这完成存在性的归纳。 对任一区间和，复合 \(V_i\to V_j\) 在一个直和项上秩为一，当且仅当该区间同时包含 \(i,j\)。故
	\[
		r_{ij}=\sum_{\substack{1\le u\le i\\j\le v\le n}}m_{uv}.
	\]
	固定 \(a\le b\)，先作
	\[
		r_{ab}-r_{a-1,b}=\sum_{v=b}^n m_{av}.
	\]
	这一步消去起点早于 \(a\) 的区间，留下恰从 \(a\) 开始、终点不早于 \(b\) 的区间。将 \(b\) 换成 \(b+1\) 得
	\[
		r_{a,b+1}-r_{a-1,b+1}=\sum_{v=b+1}^n m_{av}.
	\]
	两式再相减，消去终点晚于 \(b\) 的区间，只留下
	\[
		\bigl(r_{ab}-r_{a-1,b}\bigr) -\bigl(r_{a,b+1}-r_{a-1,b+1}\bigr)=m_{ab}.
	\]
	当 \(a=1\) 或 \(b=n\) 时，约定的边界秩为零，空和也为零。这两次消去分别固定精确起点和精确终点，展开便是所列四项差。每个区间不可分解且不同构已经证明；反过来，一个不可分解表示的区间分解只能有一项，因此类数为区间数。
\end{proof}
这个公式也给出判别同构的有限算法：计算各段复合的秩，再取四项差即可。仅计算相邻箭的秩仍不够。例如两个表示都可有维数向量 \((1,2,1)\) 及相邻秩 \(1,1\)，但复合秩分别为一和零；它们的区间分解因而不同。

\begin{proposition}\label{b5:prop:interval-hom}
	设 \(k\) 为域，\(n\ge1\) 为整数，\(Q=(1\to2\to\cdots\to n)\)，\(1\le a\le b\le n\)、\(1\le c\le d\le n\)，\(I[a,b],I[c,d]\) 为 \(k\) 上的区间表示。则
	\[
		\dim_k\operatorname{Hom}_Q\bigl(I[a,b],I[c,d]\bigr)
		=\begin{cases}
			1,&c\le a\le d\le b,\\
			0,&\text{其他情形}.
		\end{cases}
	\]
	在非零情形，一个基态射是在重合顶点 \(a,\ldots,d\) 取恒等映射，其余分量取零的态射。
\end{proposition}
\begin{proof}
	两区间不相交时，逐顶点的映射全为零。若相交，重合区间内部的恒等箭迫使全部分量为同一个标量 \(\lambda\)。若 \(a<c\)，在目标首次出现的箭 \(c-1\to c\) 上，源的箭映射为恒等而目标的箭映射从零空间出发，交换条件迫使 \(\lambda=0\)。若 \(b<d\)，在源最后出现后的箭 \(b\to b+1\) 上，源的箭映射进入零空间，而目标的箭映射为恒等，同样迫使 \(\lambda=0\)。非零态射因此要求 \(c\le a\)、\(d\le b\)，连同重合条件得到 \(c\le a\le d\le b\)。 这些条件成立时，在重合顶点 \(a,\ldots,d\) 取恒等映射，其余分量取零。重合区间内部的箭交换；从其左端进入及从其右端离开的箭，两边复合都为零；其余箭上的复合也都为零。因此得到态射，且任意态射都是它的标量倍，所以空间恰为一维。
\end{proof}
在 \(A_4=(1\to2\to3\to4)\) 上取 \(X=I[2,4]\)、\(Y=I[1,3]\)。非零态射 \(\varphi:X\to Y\) 的逐顶点分量为 \((0,1,1,0)\)。在箭 \(1\to2\) 上，两边复合都是 \(0\)；在 \(2\to3\) 上，都是 \(1\)；在 \(3\to4\) 上，又都是 \(0\)。反向态射 \(\psi:Y\to X\) 在重合顶点 \(2,3\) 的分量须同为 \(\lambda\)，但箭 \(1\to2\) 的交换条件已给出 \(\lambda\cdot1=0\cdot0=0\)，所以反向同态空间为零。这也解释了前面旗保持换基的方向：对共同结束于 \(n-1\) 的区间，非零同态 \(I[a,n-1]\to I[b,n-1]\) 恰要求 \(b\le a\)，允许较晚出现的末端方向加上较早传播来的方向。
